\begin{helper}[Product Model Monotonicity]
Fix \(M\in\mathbb N\), \(p_0\in[0,1]\), and sets \(P_R, P_B\subseteq\mathrm{Fin}(M)\).
Let \(q=\noderef{FixedSetProductModelMassFn}(p_0, P_R, P_B)\).
For any two families of assignments \(\mathcal S\subseteq\mathcal T\), we have the finite-sum monotonicity:
\[
\sum_{v\in\mathcal S}\left(\prod_i q(v_i)\right) \le \sum_{v\in\mathcal T}\left(\prod_i q(v_i)\right).
\]
\end{helper}
\begin{proof}
By \(\noderef{FixedSetProductModelMass}\), applied with the hypotheses
\(0\le p_0\le1\), the mass function \(q\) is pointwise nonnegative:
\[
q(a)\ge0\qquad\text{for every one-coordinate value }a.
\]
Therefore, for every product assignment \(v\),
\[
W(v):=\prod_{i\in\mathrm{Fin}(2M)}q(v_i)\ge0,
\]
because a finite product of nonnegative factors is nonnegative.

Since \(\mathcal S\subseteq\mathcal T\), decompose the finite set
\(\mathcal T\) as the disjoint union of \(\mathcal S\) and
\(\mathcal T\setminus\mathcal S\).  Finite additivity of sums gives
\[
\sum_{v\in\mathcal T}W(v)
=
\sum_{v\in\mathcal S}W(v)
+
\sum_{v\in\mathcal T\setminus\mathcal S}W(v).
\]
Every term in the last sum is nonnegative, so the last sum is nonnegative.
Hence
\[
\sum_{v\in\mathcal S}W(v)\le \sum_{v\in\mathcal T}W(v),
\]
which is the claimed monotonicity.
\end{proof}
